\chapter[phd_introduction]{Introduction}
 [
  A PhD manuscript uses exactly the same commands as an HDR one: only the \texttt{phd} option of the document class differs. \\
  This chapter introduces the placeholder problem addressed in this thesis, and outlines the rest of the document.
 ]

\section[phd_context]{Context}

Lorem ipsum dolor sit amet, consectetur adipiscing elit, sed do eiusmod tempor incididunt ut labore et dolore magna aliqua~\cite{ipsum1998foundational}.
Ut enim ad minim veniam, quis nostrud exercitation ullamco laboris nisi ut aliquip ex ea commodo consequat~\cite{tempor2009textbook}.

The rest of this document uses \acro{ml} vocabulary, and in particular the \acrofull{cnn}.
Note that \acro{cnn_plural} is a separate entry of \verb|acronyms.tex|, whose only purpose is to provide a plural tooltip: its full name is \acroname{cnn_plural}, and it does not appear in the list of acronyms.
Recall that \acrofull{sota} results in this placeholder field are due to \cite{amet2017seminal}.

\section[phd_problem]{Problem Statement}

\begin{boxenv}[phd_problem_box]{Problem}[Placeholder problem]
    Given a set of placeholder observations $\{\bm{x}_i\}_{i = 1}^N$ and their placeholder targets $\{\bm{y}_i\}_{i = 1}^N$, find a placeholder model $f_\theta$ that generalizes to unseen placeholder observations.
\end{boxenv}

Duis aute irure dolor in reprehenderit in voluptate velit esse cillum dolore eu fugiat nulla pariatur.
\nameref{phd_problem_box} is addressed in \nameref{phd_contribution}, after the review of \nameref{phd_state_of_the_art}.

\section[phd_outline]{Outline}

\nameref{phd_state_of_the_art} reviews the placeholder literature.
\nameref{phd_contribution} presents our contribution, evaluated in \nameref{phd_contribution_results}.
\nameref{phd_conclusion} concludes, and \nameref{phd_proofs} gathers the proofs.
